\documentclass[../../../../master.tex]{subfiles}
\begin{document}

\monographleaf{Domain 05: Transformers \& Large Language Models}{LLM Pretraining, Adaptation, Prompting \& Efficient Inference}{Prompting, Decoding, KV Caching \& Test-Time Scaling}{Prompt Engineering, Decoding Strategies, KV Caching \& Test-Time Scaling}{Foundations of LLMs Ch.~3 \& 5, Raschka Sec.~5.3, Jurafsky \& Martin Ch.~7.6, 10, 12.4, RLHF Ch.~7}

\section{Problem Formulation \& Unified Notation}\label{sec:decode:formulation}

Let $\mathcal{V}$ denote a discrete token vocabulary of size $|\mathcal{V}|$. Given an input prompt $\vecb{x} = (x_0, x_1, \dots, x_m) \in \mathcal{V}^{m+1}$ (with $x_0 = \langle\mathrm{SOS}\rangle$), an autoregressive decoder-only Transformer parameterized by $\theta$ defines a conditional probability distribution over an output sequence $\vecb{y} = (y_1, \dots, y_n) \in \mathcal{V}^n$ via the probability chain rule \bookcite{foundations_large_language_models}{Ch.~5, pp.~211--213}:
\begin{equation}\label{eq:decode:autoregressive_factorization}
  \log \Pr_{\theta}(\vecb{y} \mid \vecb{x})
  = \sum_{i=1}^{n} \log \Pr_{\theta}(y_i \mid \vecb{x}, \vecb{y}_{<i}),
  \qquad \vecb{y}_{<i} \triangleq (y_1, \dots, y_{i-1}).
\end{equation}
At decoding step $i$, let $i' = m + i$ denote the total context length of the concatenated sequence $[\vecb{x}, \vecb{y}_{<i}]$. The Transformer stack produces hidden representations $\mat{H} = \mathrm{Dec}([\vecb{x}, \vecb{y}_{<i}]) \in \R^{i' \times d}$, and the unembedding matrix $\mat{W}_o \in \R^{d \times |\mathcal{V}|}$ maps the terminal hidden vector $\vecb{h}_{m+i} \in \R^d$ to unnormalized logits $\vecb{u}_{i} = \mat{W}_o\T \vecb{h}_{m+i} \in \R^{|\mathcal{V}|}$ \bookcite{speech_language_processing_aug26}{Sec.~7.6, p.~196}. The next-token categorical distribution is
\begin{equation}\label{eq:decode:softmax_logits}
  \Pr_{\theta}(\cdot \mid \vecb{x}, \vecb{y}_{<i})
  = \softmax(\vecb{u}_i)
  = \softmax\!\left((\mat{H}\mat{W}_o)_{m+i}\right).
\end{equation}
LLM inference decomposes into two coupled sub-problems: (i) \emph{model computation}, evaluating Eq.~\eqref{eq:decode:softmax_logits} under strict latency and HBM bandwidth constraints via Key-Value (KV) caching and parallel verification; and (ii) \emph{search and conditioning}, shaping $\vecb{x}$ (in-context learning, Chain-of-Thought) and traversing or scaling the output space $\mathcal{V}^n$ (sampling, beam search, test-time verification).

\begin{crossbook}[Cross-Book Notation \& Pedagogical Reconciliation]
Source textbooks adopt distinct conventions for logits, temperature, sequence indexing, and verification models. Table~\ref{tab:decode:notation} reconciles these notations into a single unified formalism used throughout this monograph.
\vspace{4pt}

{\footnotesize
\noindent
\begin{tabularx}{\linewidth}{@{}>{\bfseries\color{monographNavy}\raggedright\arraybackslash}p{0.18\linewidth} >{\raggedright\arraybackslash}p{0.26\linewidth} >{\raggedright\arraybackslash}X@{}}
\toprule
\textbf{Unified Symbol} & \textbf{Meaning in This Monograph} & \textbf{Source Textbook Equivalents} \\
\midrule
$\vecb{x}, \vecb{y}_{<i}, y_i$ & Prompt ($m{+}1$ tokens), prefix, next token & $\vecb{x}, \vecb{y}_{<i}, y_i$ \bookcite{foundations_large_language_models}{p.~211}; $w_{<t}, w_t$ \bookcite{speech_language_processing_aug26}{p.~196}; $s_i, a_i$ \bookcite{reinforcement_learning_from_human_feedback}{p.~103} \\
$\vecb{u}_i \in \R^{|\mathcal{V}|}$ & Unnormalized next-token logit vector & $(\mat{H}\mat{W}_o)_{m+i}$ \bookcite{foundations_large_language_models}{p.~212}; $\vecb{u}$ \bookcite{speech_language_processing_aug26}{p.~196}; \texttt{logits} \bookcite{build_large_language_model_from_scratch}{p.~381} \\
$\tau > 0$ & Softmax temperature scaling parameter & $\tau$ \bookcite{speech_language_processing_aug26}{p.~198}; $T/\tau$ \bookcite{foundations_large_language_models}{p.~224}; \texttt{temperature} \bookcite{build_large_language_model_from_scratch}{p.~384} \\
$p(\cdot), q(\cdot)$ & Target (verifier) and draft distributions & $\Pr_p(\cdot), \Pr_q(\cdot)$ \bookcite{foundations_large_language_models}{p.~226} \\
$b, s, n_{\mathrm{layers}}, n_{\mathrm{kv}}, d_h$ & Batch, seq len, layers, KV heads, head dim & $L$ layers, $d$ hidden dim \bookcite{foundations_large_language_models}{p.~216} \\
\bottomrule
\end{tabularx}
}
\captionof{table}{Cross-book notation reconciliation for prompting, decoding, KV caching, and test-time scaling.}\label{tab:decode:notation}
\end{crossbook}

\section{In-Context Learning, Prompting \& Mechanistic Circuits}\label{sec:decode:prompting}

\begin{sloppypar}
Unlike gradient-based fine-tuning, \emph{in-context learning} (ICL) adapts model behavior during the forward pass at inference time with zero parameter updates ($\Delta\theta = \zeros$, see \bookcite{foundations_large_language_models}{Ch.~3, p.~106}). Prompting spans a spectrum of conditioning regimes:
\begin{itemize}
  \item \textbf{Zero-Shot, One-Shot \& Few-Shot ICL:} In zero-shot prompting, the input $\vecb{x}$ specifies task instructions and role constraints. One-shot and few-shot prompting prepend $k$ input--output demonstrations $\{(x^{(j)}, y^{(j)})\}_{j=1}^k$ to $\vecb{x}$, inducing an empirical task mapping directly in the residual stream \bookcite{foundations_large_language_models}{Ch.~3, pp.~106--108}.
  \item \textbf{Chain-of-Thought (CoT) \& Problem Decomposition:} Rather than mapping $\vecb{x} \mapsto \vecb{y}_{\mathrm{ans}}$ directly, CoT conditions generation on an explicit intermediate reasoning trajectory $\vecb{z} = (z_1, \dots, z_T)$ before emitting the final answer: $\Pr(\vecb{y}_{\mathrm{ans}} \mid \vecb{x}) = \sum_{\vecb{z}} \Pr(\vecb{y}_{\mathrm{ans}} \mid \vecb{x}, \vecb{z}) \Pr(\vecb{z} \mid \vecb{x})$ \bookcite{foundations_large_language_models}{Ch.~3, pp.~122--125}. Decomposing multi-step problems allocates serial Transformer compute across intermediate token positions.
  \item \textbf{Self-Consistency \& Tree-of-Thoughts (ToT):} Sampling $M$ diverse CoT paths $\{\vecb{z}^{(m)}, \vecb{y}_{\mathrm{ans}}^{(m)}\}_{m=1}^M$ and marginalizing via majority vote approximates the intractable sum over $\vecb{z}$, whereas ToT structures intermediate thoughts as a search tree with lookahead and backtracking (\bookcite{foundations_large_language_models}{Ch.~3, pp.~125--141}).
\end{itemize}
\end{sloppypar}

\begin{intuition}[Mechanistic Interpretability: Logit Lens, SAEs \& Induction Heads]
How does a frozen Transformer execute ICL without weight updates? Mechanistic interpretability probes the residual stream across layers $l \in \{1, \dots, L\}$ \bookcite{speech_language_processing_aug26}{Ch.~10, pp.~248--251}:
\begin{enumerate}
  \item \textbf{The Logit Lens:} Projecting intermediate hidden states $\vecb{h}_t^{(l)}$ directly through the unembedding matrix $\softmax(\mat{W}_o\T \vecb{h}_t^{(l)})$ reveals a layer-wise refinement workspace where early layers resolve entity attributes and later layers format the target language output.
  \item \textbf{Sparse Autoencoders (SAEs):} Decompose polysemantic neuron activations into sparse, monosemantic feature dictionaries for causal intervention.
  \item \textbf{Induction Heads:} Two-layer attention circuits implementing the pattern-completion rule $A\,B \dots A \to B$ (and its fuzzy semantic generalization $A^*\,B^* \dots A \to B$). When attending from the current token $A$, a previous-token head writes $A$'s identity into $B$'s key, allowing the induction head to perform \emph{prefix matching} back to the prior $A$ and \emph{copy} the subsequent token $B$ (Fig.~\ref{fig:decode:induction_head}). Ablating induction heads causally degrades few-shot ICL performance \bookcite{speech_language_processing_aug26}{Ch.~10, pp.~250--251}.
\end{enumerate}
\end{intuition}

\monographfig[0.88\textwidth]{speech_language_processing_aug26_p259_fig1.png}{Induction head circuit implementing in-context pattern completion: attending from the second occurrence of \emph{vintage}, the prefix-matching mechanism locates the prior occurrence and the copying mechanism boosts the probability of \emph{cars}. Source: \bookcite{speech_language_processing_aug26}{Ch.~10, Fig.~10.7, p.~251 (PDF p.~259)}.}{fig:decode:induction_head}

\section{Core Mathematical Derivations: Decoding \& Sampling}\label{sec:decode:decoding}

Exact maximum a posteriori (MAP) decoding $\argmax_{\vecb{y}} \log \Pr_{\theta}(\vecb{y} \mid \vecb{x})$ over a search tree of depth $n_{\max}$ and branching factor $|\mathcal{V}|$ requires evaluating $|\mathcal{V}|^{n_{\max}}$ sequences, which is computationally intractable \bookcite{foundations_large_language_models}{Ch.~5, pp.~219--220}. Practical decoding algorithms prune the level-$i$ candidate set via $\mathcal{Y}_i = \mathrm{Prune}(\mathcal{Y}_{i-1} \times \mathcal{V})$.

\begin{coremath}[Deterministic Search: Greedy Decoding \& Beam Search]
At step $i$, the cumulative log-probability of extending prefix $\vecb{y}_{<i}$ with token $y_i$ decomposes additively:
\begin{equation}\label{eq:decode:cumulative_logprob}
  \log \Pr(\vecb{y}_{1:i} \mid \vecb{x}) = \log \Pr(\vecb{y}_{<i} \mid \vecb{x}) + \log \Pr(y_i \mid \vecb{x}, \vecb{y}_{<i}).
\end{equation}
\begin{itemize}
  \item \textbf{Greedy Decoding:} Sets $|\mathcal{Y}_i| = 1$ by selecting $\hat{y}_i = \argmax_{y_i \in \mathcal{V}} \log \Pr(y_i \mid \vecb{x}, \vecb{y}_{<i})$ \bookcite{speech_language_processing_aug26}{Sec.~7.6.1, p.~196}. Greedy decoding is $O(n|\mathcal{V}|)$ but is locally myopic and prone to repetitive degeneration.
  \item \textbf{Beam Search:} Maintains the top-$K$ hypotheses $\mathcal{Y}_i = \mathrm{argTopK}_{\vecb{y} \in \mathcal{Y}_{i-1} \times \mathcal{V}} \log \Pr(\vecb{y} \mid \vecb{x})$ with beam width $K$ (typically $K \in \{2, 4\}$) \bookcite{foundations_large_language_models}{Ch.~5, pp.~221--222}. Beam search excels in constrained sequence-to-sequence tasks (e.g., machine translation) but often produces generic, low-entropy text in open-ended generation.
\end{itemize}
\end{coremath}

To restore diversity and match human-like entropy profiles, stochastic decoding samples $y_i \sim p_{\tau}(\cdot \mid \vecb{x}, \vecb{y}_{<i})$ after shaping and truncating the logit distribution.

\begin{coremath}[Temperature Scaling, Top-$k$ \& Top-$p$ (Nucleus) Filtering]
Given logits $\vecb{u} \in \R^{|\mathcal{V}|}$ and temperature $\tau > 0$, the temperature-scaled distribution is \bookcite{build_large_language_model_from_scratch}{Sec.~5.3.1, pp.~380--387}, \bookcite{speech_language_processing_aug26}{Sec.~7.6.3, p.~198}:
\begin{equation}\label{eq:decode:temperature_softmax}
  p_{\tau}(w_j) = \frac{\exp(u_j / \tau)}{\sum_{v=1}^{|\mathcal{V}|} \exp(u_v / \tau)}, \qquad
  \lim_{\tau \to 0^+} p_{\tau}(w_j) = \mathbb{I}[j = \argmax_v u_v], \quad
  \lim_{\tau \to \infty} p_{\tau}(w_j) = \frac{1}{|\mathcal{V}|}.
\end{equation}
Sampling directly from $p_{\tau}$ across all $|\mathcal{V}|$ tokens exposes generation to the \emph{unreliable tail} (accumulated probability mass over thousands of implausible tokens). Two truncation filters eliminate this tail:
\begin{enumerate}
  \item \textbf{Top-$k$ Filtering:} Retain the $k$ highest-logit indices $\mathcal{V}_k = \mathrm{argTopK}(\vecb{u}, k)$ and mask all other logits to $-\infty$ prior to softmax \bookcite{build_large_language_model_from_scratch}{Sec.~5.3.2, pp.~388--393}:
  \begin{equation}\label{eq:decode:topk}
    \tilde{u}_j = \begin{cases} u_j & \text{if } j \in \mathcal{V}_k, \\ -\infty & \text{otherwise}, \end{cases}
    \qquad p_{\tau, k}(w_j) = \frac{\exp(\tilde{u}_j / \tau)}{\sum_{v \in \mathcal{V}_k} \exp(u_v / \tau)}.
  \end{equation}
  \item \textbf{Top-$p$ (Nucleus) Filtering:} Sort tokens in descending probability order $p_{\tau}(w_{(1)}) \ge p_{\tau}(w_{(2)}) \ge \dots \ge p_{\tau}(w_{(|\mathcal{V}|)})$ and define the minimal nucleus $\mathcal{V}_p = \{w_{(1)}, \dots, w_{(m^*)}\}$ whose cumulative mass meets threshold $p \in (0, 1]$ \bookcite{foundations_large_language_models}{Ch.~5, pp.~222--224}, \bookcite{speech_language_processing_aug26}{Sec.~7.6.4, p.~200}:
  \begin{equation}\label{eq:decode:topp}
    m^* = \min \left\{ m \in \{1, \dots, |\mathcal{V}|\} : \sum_{r=1}^{m} p_{\tau}(w_{(r)}) \ge p \right\},
    \qquad p_{\tau, p}(w_j) = \frac{p_{\tau}(w_j)\,\mathbb{I}[w_j \in \mathcal{V}_p]}{\sum_{v \in \mathcal{V}_p} p_{\tau}(v)}.
  \end{equation}
\end{enumerate}
\end{coremath}

\begin{intuition}[Why Nucleus ($p$) Adapts Better Than Fixed $k$]
When the model is highly certain (e.g., completing a deterministic idiom), $p_{\tau}(w_{(1)}) \approx 0.95$, so a fixed $k = 50$ forces 49 junk tokens into the candidate pool, whereas top-$p$ with $p = 0.9$ automatically collapses $|\mathcal{V}_p| = 1$. Conversely, in high-entropy contexts (e.g., starting a story), probability mass spreads across hundreds of valid continuations; top-$p$ expands $|\mathcal{V}_p|$ dynamically while still excising the long tail.
\end{intuition}

\begin{algorithm}[htbp]
\caption{Temperature-Scaled Top-$k$ / Top-$p$ Nucleus Sampling}\label{alg:decode:nucleus}
\KwIn{Context $\vecb{x}$, model $\theta$, temperature $\tau \ge 0$, top-$k$ bound $k \le |\mathcal{V}|$, nucleus $p \in (0, 1]$, max tokens $n_{\max}$}
\KwOut{Generated sequence $\vecb{y}$}
Initialize $\vecb{y} \leftarrow ()$, $\mathrm{cache} \leftarrow \mathrm{Dec}_{\mathrm{kv}}(\vecb{x})$ \tcp*{Prefill phase}
\For{$i \leftarrow 1$ \KwTo $n_{\max}$}{
  Compute logits $\vecb{u}_i \in \R^{|\mathcal{V}|}$ using $\mathrm{cache}$\;
  \If{$\tau = 0$}{
    $y_i \leftarrow \argmax_{v \in \mathcal{V}} u_{i,v}$ \tcp*{Greedy fallback}
  }
  \Else{
    \If{$k < |\mathcal{V}|$}{
      $\mathcal{V}_k \leftarrow \mathrm{argTopK}(\vecb{u}_i, k)$;\quad set $u_{i,v} \leftarrow -\infty$ for $v \notin \mathcal{V}_k$\;
    }
    $\vecb{p} \leftarrow \softmax(\vecb{u}_i / \tau)$\;
    Sort $\vecb{p}$ descending: $p_{(1)} \ge p_{(2)} \ge \dots$; find $m^* = \min \{m : \sum_{r=1}^m p_{(r)} \ge p\}$\;
    $\mathcal{V}_p \leftarrow \{(1), \dots, (m^*)\}$;\quad zero out $p_v$ for $v \notin \mathcal{V}_p$ and renormalize $\vecb{p} \leftarrow \vecb{p} / \sum_{v \in \mathcal{V}_p} p_v$\;
    Sample $y_i \sim \Cat(\vecb{p})$\;
  }
  Append $y_i$ to $\vecb{y}$ and update $\mathrm{cache}$\;
  \If{$y_i = \langle\mathrm{EOS}\rangle$}{\textbf{break}\;}
}
\Return{$\vecb{y}$}
\end{algorithm}

\section{KV Caching, Prefill vs. Decode \& Efficient Serving}\label{sec:decode:kv_cache}

In causal self-attention, keys $\vecb{k}_{t}$ and values $\vecb{v}_{t}$ depend only on tokens $\le t$ and never change in subsequent steps. Recomputing $\mat{K}_{1:i'}, \mat{V}_{1:i'}$ from scratch at every step incurs $O(n^2)$ redundant work. Storing past key--value pairs in a \emph{KV cache} $\mathrm{cache} = \{(\mat{K}^{(l)}, \mat{V}^{(l)})\}_{l=1}^{n_{\mathrm{layers}}}$ splits inference into two hardware-distinct phases \bookcite{foundations_large_language_models}{Ch.~5, pp.~214--217} (Fig.~\ref{fig:decode:prefill_decode_book}):
\begin{enumerate}
  \item \textbf{Prefilling (Compute-Bound):} Processes all $m$ prompt tokens in parallel via $\mathrm{cache} = \mathrm{Dec}_{\mathrm{kv}}(\vecb{x})$, executing dense GEMMs with high arithmetic intensity and dominating Time-to-First-Token (TTFT).
  \item \textbf{Decoding (Memory-Bandwidth-Bound):} Generates one token at a time, appending $(\vecb{k}_{i'}, \vecb{v}_{i'})$ to $\mathrm{cache}$ and attending $\vecb{q}_{i'}$ against all $i'$ cached keys/values, governing Inter-Token Latency (ITL).
\end{enumerate}

\monographfig[0.70\textwidth]{foundations_large_language_models_p217_fig1.png}{Two-phase LLM inference: (a) Prefilling processes all $m$ input tokens in parallel to populate the initial KV cache; (b) Decoding generates one token per step, appending $(\vecb{k}_{m+n}, \vecb{v}_{m+n})$ to the KV cache and attending $\vecb{q}_{m+n}$ across all prior keys and values. Source: \bookcite{foundations_large_language_models}{Ch.~5, Fig.~5.3, p.~217}.}{fig:decode:prefill_decode_book}

\begin{coremath}[KV Cache Memory Footprint \& Arithmetic Intensity]
For a Transformer with $n_{\mathrm{layers}}$ layers, $n_{\mathrm{kv}}$ key-value heads per layer, head dimension $d_h$, batch size $b$, total sequence length $s = m + n$, and $p_{\mathrm{bytes}}$ bytes per element (e.g., $p_{\mathrm{bytes}} = 2$ for FP16/BF16), the total KV cache footprint is
\begin{equation}\label{eq:decode:kv_memory}
  M_{\mathrm{KV}} = 2 \cdot b \cdot s \cdot n_{\mathrm{layers}} \cdot n_{\mathrm{kv}} \cdot d_h \cdot p_{\mathrm{bytes}} \quad [\text{bytes}],
\end{equation}
where the factor of $2$ accounts for storing both $\mat{K}$ and $\mat{V}$ tensors. Architectural variants reduce $n_{\mathrm{kv}}$ relative to query heads $n_q$:
\begin{itemize}
  \item \textbf{Multi-Head Attention (MHA):} $n_{\mathrm{kv}} = n_q$ (baseline $1\times$).
  \item \textbf{Grouped-Query Attention (GQA):} $1 < n_{\mathrm{kv}} < n_q$, sharing each KV head across $n_q / n_{\mathrm{kv}}$ query heads (e.g., $8\times$ memory reduction for $n_q=64, n_{\mathrm{kv}}=8$).
  \item \textbf{Multi-Query Attention (MQA):} $n_{\mathrm{kv}} = 1$, achieving maximal $n_q\times$ KV reduction.
\end{itemize}
\textbf{Arithmetic Intensity Gap:} During a single decode step over $b$ sequences, each linear projection performs $\approx 2 b P_{\mathrm{layer}}$ FLOPs while loading $\approx P_{\mathrm{layer}} \cdot p_{\mathrm{bytes}}$ weight bytes plus $M_{\mathrm{KV}} / n_{\mathrm{layers}}$ cache bytes from HBM. For small $b$, arithmetic intensity $\mathcal{I}_{\mathrm{decode}} = \frac{\mathrm{FLOPs}}{\mathrm{Bytes}} \approx \frac{2b}{p_{\mathrm{bytes}}} \sim O(b)$ sits far below modern GPU ridge points ($\sim 200$--$300\;\mathrm{FLOPs/byte}$), making decoding strictly memory-bandwidth-bound.
\end{coremath}

To overcome HBM capacity and bandwidth bottlenecks in production serving, modern inference engines combine four complementary techniques \bookcite{foundations_large_language_models}{Ch.~5, pp.~230--241}:
\begin{itemize}
  \item \textbf{PagedAttention:} Partitions each sequence's KV cache into fixed-size non-contiguous physical blocks managed via a block table (analogous to OS virtual memory), eliminating external fragmentation and enabling copy-on-write prefix sharing \bookcite{foundations_large_language_models}{Ch.~5, pp.~237--238}.
  \item \textbf{Continuous Batching \& Chunked Prefilling:} Iteration-level scheduling injects new requests into active batches as soon as sequences emit $\langle\mathrm{EOS}\rangle$, while chunked prefilling splits long prompts into fixed-size chunks overlapped with decode steps to bound ITL variance \bookcite{foundations_large_language_models}{Ch.~5, pp.~235--239}.
  \item \textbf{IO-Aware Tiling (FlashAttention):} Tiles $\mat{Q}, \mat{K}, \mat{V}$ into GPU SRAM and maintains running $(\max, \text{sum})$ statistics for online softmax, avoiding $O(s^2)$ HBM materialization of the attention matrix.
  \item \textbf{KV \& Weight Quantization:} Compresses weights (e.g., INT4/FP8) and KV caches (INT8/FP8) to halve or quarter HBM traffic per decode step, trading minor calibration error for proportional TPS gains \bookcite{foundations_large_language_models}{Ch.~5, pp.~240--241}.
\end{itemize}

\begin{figure}[htbp]
\centering
\begin{tikzpicture}[node distance=0.7cm and 0.7cm, >=Stealth]
  % Left: Prefill
  \node[monoBox, minimum width=2.4cm] (prompt) {\textbf{Prompt} $\vecb{x}_{0:m}$\\($m$ tokens)};
  \node[monoTealBox, below=0.7cm of prompt, minimum width=2.4cm] (prefill) {\textbf{Prefill Phase}\\$O(m^2)$ Compute-Bound\\Parallel GEMMs};
  \node[monoIndigoBox, below=0.7cm of prefill, minimum width=2.4cm] (kv) {\textbf{KV Cache}\\$2 b s n_{\mathrm{layers}} n_{\mathrm{kv}} d_h p_{\mathrm{bytes}}$\\(Paged Blocks)};

  % Middle: Autoregressive Decode
  \node[monoBox, right=0.75cm of prefill, minimum width=2.65cm] (decode) {\textbf{Autoregressive Decode}\\$O(s)$ Memory-BW-Bound\\1 token/step};
  \node[monoTealBox, below=0.7cm of decode, minimum width=2.65cm] (sample) {\textbf{Logit Filter}\\Temp $\tau$, Top-$k$, Top-$p$};

  % Right: Speculative Decoding
  \node[monoIndigoBox, right=0.75cm of decode, minimum width=2.9cm] (draft) {\textbf{Draft Model} $q(\cdot)$\\Generate $\tau$ tokens $\hat{y}_{i+1:i+\tau}$\\(Fast, Small)};
  \node[monoTealBox, below=0.7cm of draft, minimum width=2.9cm] (verify) {\textbf{Target Verify} $p(\cdot)$\\1 Parallel Forward Pass\\for all $\tau{+}1$ positions};
  \node[monoWarnBox, below=0.7cm of verify, minimum width=2.9cm] (reject) {\textbf{Lossless Rejection}\\Accept if $r_t \le p/q$;\\else resample $p_{\mathrm{res}}$};

  % Arrows
  \draw[monoArrow] (prompt) -- (prefill);
  \draw[monoArrow] (prefill) -- (kv);
  \draw[monoArrow] (kv.east) -- ++(0.35,0) |- (decode.west);
  \draw[monoArrow] (decode) -- (sample);
  \draw[monoArrow] (sample.west) -- node[midway, below=1pt, font=\scriptsize, color=monographSlate] {append} (kv.east);

  \draw[monoArrow] (draft) -- (verify);
  \draw[monoArrow] (verify) -- (reject);
  \draw[monoArrow, dashed] (reject.east) -- ++(0.35,0) |- node[pos=0.25, right, font=\scriptsize, color=monographSlate] {accept $n_a{+}1$} (draft.east);
\end{tikzpicture}
\caption{Unified inference architecture: (Left) Compute-bound prefill populates the paged KV cache; (Center) Memory-bandwidth-bound autoregressive decode appends one token per step; (Right) Speculative decoding amortizes target-model HBM reads by verifying $\tau$ draft tokens in a single parallel pass.}\label{fig:decode:prefill_kv_speculative}
\end{figure}

Figure~\ref{fig:decode:kv_memory_scaling} quantifies the impact of Eq.~\eqref{eq:decode:kv_memory} for a 70B-class Transformer ($n_{\mathrm{layers}}=80, d_h=128$, FP16 $p_{\mathrm{bytes}}=2$, single request $b=1$) across MHA ($n_{\mathrm{kv}}=64$), GQA-8 ($n_{\mathrm{kv}}=8$), GQA-4 ($n_{\mathrm{kv}}=4$), and MQA ($n_{\mathrm{kv}}=1$). At $s=128\mathrm{k}$ tokens, MHA consumes $320\;\mathrm{GiB}$ for a single sequence's KV cache alone, whereas GQA-8 reduces this to $40\;\mathrm{GiB}$ and MQA to $5\;\mathrm{GiB}$.

\begin{figure}[htbp]
\centering
\begin{tikzpicture}
\begin{axis}[
  width=0.86\textwidth,
  height=6.2cm,
  xlabel={Total Context Length $s$ (thousands of tokens)},
  ylabel={KV Cache Memory per Sequence (GiB, FP16)},
  xmin=1, xmax=128,
  ymin=0.01, ymax=500,
  xmode=log,
  ymode=log,
  log basis x={2},
  log basis y={10},
  xtick={1,2,4,8,16,32,64,128},
  xticklabels={1k,2k,4k,8k,16k,32k,64k,128k},
  grid=both,
  major grid style={color=monographRule, line width=0.4pt},
  minor grid style={color=monographRule!40, line width=0.2pt},
  legend pos=south east,
  legend style={font=\small, draw=monographRule, fill=white, fill opacity=0.9},
  tick label style={font=\small, color=monographSlate},
  label style={font=\small\bfseries, color=monographNavy}
]
% M = 2 * 1 * (s*1024) * 80 * n_kv * 128 * 2 / (1024^3) = s * n_kv * (40960 / 1048576) = s * n_kv * 0.0390625 * 0.5?
% Wait: 2 * 1024 * 80 * 128 * 2 = 41,943,040 bytes per 1k tokens per kv head.
% 41,943,040 / (1024^3) = 0.0390625 / 1? Let's check: 41943040 / 1073741824 = 0.0390625 GB per 1k tokens per KV head!
% Wait: 2 * 80 * 128 * 2 = 40960 bytes/token/kv_head. * 1024 = 41,943,040 bytes = 40 MiB = 0.0390625 GiB.
% For MHA (n_kv=64): 64 * 0.0390625 = 2.5 GiB per 1k tokens -> at 128k: 320 GiB!
% Wait: let's set ymax=400 so 320 fits cleanly!
\addplot[color=monographCrimson, thick, mark=square*, mark size=1.8pt, domain=1:128, samples=8] {2.5 * x};
\addlegendentry{MHA ($n_{\mathrm{kv}} = 64$, $2.50\;\mathrm{GiB/1k}$)}

\addplot[color=monographNavy, thick, mark=*, mark size=1.8pt, domain=1:128, samples=8] {0.3125 * x};
\addlegendentry{GQA-8 ($n_{\mathrm{kv}} = 8$, $0.31\;\mathrm{GiB/1k}$)}

\addplot[color=monographTeal, thick, mark=triangle*, mark size=2pt, domain=1:128, samples=8] {0.15625 * x};
\addlegendentry{GQA-4 ($n_{\mathrm{kv}} = 4$, $0.16\;\mathrm{GiB/1k}$)}

\addplot[color=monographIndigo, thick, dashed, mark=diamond*, mark size=2pt, domain=1:128, samples=8] {0.0390625 * x};
\addlegendentry{MQA ($n_{\mathrm{kv}} = 1$, $0.039\;\mathrm{GiB/1k}$)}

% 80GB H100 VRAM reference line
\addplot[color=monographSlate, dotted, thick, domain=1:128] {80};
\node[anchor=south west, font=\scriptsize, color=monographSlate] at (axis cs:1.1,82) {Single 80\,GiB GPU VRAM};
\end{axis}
\end{tikzpicture}
\caption{KV cache memory scaling (Eq.~\eqref{eq:decode:kv_memory}) for a 70B-class Transformer ($n_{\mathrm{layers}}=80$, $d_h=128$, FP16, $b=1$). Under MHA ($n_{\mathrm{kv}}=64$), a single $128\mathrm{k}$-token context requires $320\;\mathrm{GiB}$, exceeding a single $80\;\mathrm{GiB}$ GPU by $4\times$, whereas GQA-8 reduces this to $40\;\mathrm{GiB}$ and MQA to $5\;\mathrm{GiB}$.}\label{fig:decode:kv_memory_scaling}
\end{figure}

\section{Speculative Decoding: Lossless Rejection Sampling}\label{sec:decode:speculative}

Because autoregressive decoding is memory-bandwidth-bound, evaluating $\tau + 1$ tokens in parallel on the target model $p(\cdot)$ takes nearly the same wall-clock time as evaluating a single token. Speculative decoding exploits this by using a small, fast draft model $q(\cdot)$ to autoregressively propose $\tau$ candidate tokens $\hat{y}_{i+1}, \dots, \hat{y}_{i+\tau}$, which $p(\cdot)$ verifies in one parallel forward pass \bookcite{foundations_large_language_models}{Ch.~5, pp.~225--228}.

\begin{coremath}[Proof of Exact Distribution Equivalence (Lossless Speculative Sampling)]
Let $q(y)$ and $p(y)$ denote the draft and target next-token distributions at position $t$ given prefix $(\vecb{x}, \vecb{y}_{<i}, \hat{\vecb{y}}_{1:t-1})$. Sample $\hat{y} \sim q(\cdot)$ and $r \sim \Uniform(0, 1)$:
\begin{itemize}
  \item If $r \le \frac{p(\hat{y})}{q(\hat{y})}$ (equivalently, always accept when $q(\hat{y}) \le p(\hat{y})$), \textbf{accept} $\hat{y}$.
  \item Otherwise, \textbf{reject} $\hat{y}$ and all subsequent draft tokens, and resample $\bar{y}$ from the normalized residual distribution
  \begin{equation}\label{eq:decode:residual_dist}
    p_{\mathrm{res}}(y) = \frac{\max(0,\, p(y) - q(y))}{Z}, \qquad Z = \sum_{v \in \mathcal{V}} \max(0,\, p(v) - q(v)) = 1 - \sum_{v \in \mathcal{V}} \min(p(v), q(v)).
  \end{equation}
\end{itemize}
\textbf{Theorem (Losslessness):} The emitted token $Y$ satisfies $\Pr(Y = y) = p(y)$ for all $y \in \mathcal{V}$.

\emph{Proof.} By the law of total probability, the overall acceptance probability is $\alpha = \sum_{v \in \mathcal{V}} q(v) \min\!\left(1, \frac{p(v)}{q(v)}\right) = \sum_{v \in \mathcal{V}} \min(p(v), q(v)) = 1 - Z$. For any $y \in \mathcal{V}$:
\begin{align}
  \Pr(Y = y) &= \Pr(\text{accept } y) + \Pr(\text{reject}) \cdot p_{\mathrm{res}}(y) \nonumber \\
  &= q(y)\min\!\left(1, \frac{p(y)}{q(y)}\right) + Z \cdot \frac{\max(0,\, p(y) - q(y))}{Z} \nonumber \\
  &= \min(p(y), q(y)) + \max(0,\, p(y) - q(y)) = p(y). \qquad \square \label{eq:decode:lossless_proof}
\end{align}
If all $\tau$ draft tokens are accepted, the target model emits one bonus token from $p(\cdot \mid \vecb{x}, \vecb{y}_{<i}, \hat{\vecb{y}}_{1:\tau})$ at zero extra forward-pass cost, yielding up to $\tau + 1$ exact target tokens per verification step.
\end{coremath}

\begin{algorithm}[htbp]
\caption{Lossless Speculative Decoding (Draft--Verify--Resample)}\label{alg:decode:speculative}
\KwIn{Prefix $(\vecb{x}, \vecb{y}_{<i})$, draft model $q$, target model $p$, lookahead $\tau \ge 1$}
\KwOut{Accepted tokens $(\tilde{y}_{i+1}, \dots, \tilde{y}_{i+n_a+1})$ distributed identically to $p$}
\For{$t \leftarrow 1$ \KwTo $\tau$}{
  Sample $\hat{y}_{i+t} \sim q(\cdot \mid \vecb{x}, \vecb{y}_{<i}, \hat{y}_{i+1:i+t-1})$\;
}
Evaluate $p(\cdot \mid \vecb{x}, \vecb{y}_{<i}, \hat{y}_{i+1:i+t-1})$ in parallel for $t = 1, \dots, \tau + 1$\;
Initialize $n_a \leftarrow \tau$\;
\For{$t \leftarrow 1$ \KwTo $\tau$}{
  Draw $r_t \sim \Uniform(0, 1)$\;
  \If{$r_t > \frac{p(\hat{y}_{i+t})}{q(\hat{y}_{i+t})}$}{
    $n_a \leftarrow t - 1$\;
    Sample $\bar{y}_{i+n_a+1} \sim p_{\mathrm{res}}(\cdot)$ via Eq.~\eqref{eq:decode:residual_dist}\;
    \textbf{break}\;
  }
}
\If{$n_a = \tau$}{
  Sample bonus token $\bar{y}_{i+\tau+1} \sim p(\cdot \mid \vecb{x}, \vecb{y}_{<i}, \hat{y}_{i+1:i+\tau})$\;
}
\Return{$(\hat{y}_{i+1}, \dots, \hat{y}_{i+n_a}, \bar{y}_{i+n_a+1})$}
\end{algorithm}

\section{Test-Time Scaling, Verifiers \& RLVR Reasoning}\label{sec:decode:test_time}

\begin{sloppypar}
Beyond pretraining and fine-tuning scaling laws, \emph{inference-time (test-time) scaling} improves accuracy on complex reasoning tasks by allocating additional compute during generation without updating parameters $\theta$ (\bookcite{foundations_large_language_models}{Ch.~5, pp.~241--251}; \bookcite{reinforcement_learning_from_human_feedback}{Ch.~7, pp.~101--105}). Test-time compute scales along two primary axes:
\end{sloppypar}
\begin{itemize}
  \item \textbf{Parallel Scaling (Best-of-$N$ \& Majority Voting):} Sample $N$ independent rollouts $\{\vecb{y}^{(1)}, \dots, \vecb{y}^{(N)}\}$ and aggregate via unweighted majority vote (Self-Consistency) or select $\argmax_{m} V(\vecb{x}, \vecb{y}^{(m)})$ using a learned verifier $V$ \bookcite{foundations_large_language_models}{Ch.~5, pp.~244--248}.
  \item \textbf{Sequential Scaling \& Step-Level Search:} Extend a single reasoning trajectory via self-refinement or explore partial reasoning states $\vecb{a}_{\le i} = (a_1, \dots, a_i)$ with a \emph{Process Reward Model} (PRM), which scores intermediate steps $V(\vecb{x}, \vecb{a}_{\le i})$ rather than only terminal outcomes (Outcome Reward Models, ORMs), enabling beam search or MCTS over thought steps \bookcite{foundations_large_language_models}{Ch.~5, pp.~248--251}.
  \item \textbf{RLVR-Elicited Long Thinking:} Reinforcement Learning with Verifiable Rewards (RLVR) trains policies $\pi_{\theta}$ against deterministic binary verifiers $r \in \{0, \gamma\}$ (e.g., exact math extraction via \verb|\boxed{}| or unit-test execution, Fig.~\ref{fig:decode:rlvr_loop}), teaching models to allocate long CoT budgets, backtrack, and self-verify natively at test time \bookcite{reinforcement_learning_from_human_feedback}{Ch.~7, pp.~101--112}.
\end{itemize}

\monographfig[0.72\textwidth]{reinforcement_learning_from_human_feedback_p103_fig1.png}{Reinforcement Learning with Verifiable Rewards (RLVR) loop: instead of a learned subjective reward model, a deterministic verifier scores completions $a_i$ with $r_i \in \{0, \gamma\}$, eliciting extended test-time reasoning chains. Source: \bookcite{reinforcement_learning_from_human_feedback}{Ch.~7, Fig.~28, p.~97 (PDF p.~103)}.}{fig:decode:rlvr_loop}

Figure~\ref{fig:decode:test_time_scaling} illustrates the empirical scaling regimes of test-time compute: naive majority voting saturates early as $N$ grows, whereas PRM-guided Best-of-$N$ and RLVR sequential deliberation continue converting inference FLOPs into higher pass@1 accuracy.

\begin{figure}[htbp]
\centering
\begin{tikzpicture}
\begin{axis}[
  width=0.84\textwidth,
  height=5.8cm,
  xlabel={Test-Time Compute Budget (Relative Samples / Thinking Tokens $N$)},
  ylabel={Reasoning Accuracy (pass@1, \%)},
  xmin=1, xmax=256,
  ymin=30, ymax=95,
  xmode=log,
  log basis x={2},
  xtick={1,2,4,8,16,32,64,128,256},
  xticklabels={1,2,4,8,16,32,64,128,256},
  grid=both,
  major grid style={color=monographRule, line width=0.4pt},
  legend pos=south east,
  legend style={font=\small, draw=monographRule, fill=white, fill opacity=0.9},
  tick label style={font=\small, color=monographSlate},
  label style={font=\small\bfseries, color=monographNavy}
]
% Greedy baseline
\addplot[color=monographSlate, dashed, thick, domain=1:256] {38};
\addlegendentry{Single-Shot Greedy ($N=1$)}

% Majority Voting (saturates)
\addplot[color=monographTeal, thick, mark=triangle*, mark size=2pt] coordinates {
  (1,38) (2,44) (4,51) (8,56) (16,60) (32,62.5) (64,64) (128,64.8) (256,65.2)
};
\addlegendentry{Parallel Majority Voting (Self-Consistency)}

% Best-of-N with PRM
\addplot[color=monographIndigo, thick, mark=*, mark size=1.8pt] coordinates {
  (1,38) (2,47) (4,56) (8,64) (16,70) (32,75) (64,79) (128,82) (256,84.5)
};
\addlegendentry{Best-of-$N$ + Process Reward Model (PRM)}

% RLVR Sequential Thinking
\addplot[color=monographNavy, ultra thick, mark=square*, mark size=1.8pt] coordinates {
  (1,42) (2,51) (4,60) (8,68) (16,75) (32,81) (64,85.5) (128,89) (256,91.5)
};
\addlegendentry{RLVR Sequential CoT + Verification}
\end{axis}
\end{tikzpicture}
\caption{Conceptual test-time compute scaling laws across inference budgets $N$ (illustrative log-linear trends synthesized from \bookcite{foundations_large_language_models}{Ch.~5, pp.~242--251} and \bookcite{reinforcement_learning_from_human_feedback}{Ch.~7, pp.~101--111}). Majority voting plateaus under correlated errors, while PRM-guided search and RLVR sequential deliberation sustain monotonic gains.}\label{fig:decode:test_time_scaling}
\end{figure}

\begin{table}[htbp]
\centering
\small
\begin{tabularx}{\textwidth}{@{}>{\bfseries\color{monographNavy}}l l l X@{}}
\toprule
\textbf{Mechanism} & \textbf{Compute / Step} & \textbf{Memory / Req.} & \textbf{Bottleneck \& Primary Trade-off} \\
\midrule
Prefill (no cache) & $O(m^2 d + m d^2)$ & $O(b m n_{\mathrm{layers}} n_{\mathrm{kv}} d_h)$ & Compute-bound; sets TTFT \\
Naive Decode (no cache) & $O(s^2 d + s d^2)$ & $O(1)$ cache & $O(n^3)$ total compute; unusable for long $s$ \\
KV-Cached Decode & $O(s d + d^2)$ & $O(b s n_{\mathrm{layers}} n_{\mathrm{kv}} d_h)$ & HBM bandwidth-bound; sets ITL / TPS \\
GQA / MQA & $O(s d + d^2)$ & $\frac{n_{\mathrm{kv}}}{n_q} \times$ MHA & Cuts KV HBM traffic $n_q / n_{\mathrm{kv}}\times$ \\
PagedAttention & $O(s d + d^2)$ & $\sim 0\%$ internal waste & Eliminates fragmentation; enables prefix sharing \\
Speculative Decoding & $O(\tau \cdot C_q + C_p)$ & $+ M_{\mathrm{KV}}(q)$ & Lossless $\sim 2$--$3\times$ speedup if $\alpha$ is high \\
Test-Time Search ($N$) & $O(N \cdot C_{\mathrm{seq}})$ & $O(N \cdot M_{\mathrm{KV}})$ & Trades $N\times$ inference FLOPs for pass@1 \\
\bottomrule
\end{tabularx}
\caption{Complexity and hardware bottleneck synthesis across LLM inference mechanisms.}\label{tab:decode:complexity}
\end{table}

\section{Numerical Pitfalls, Edge Cases \& Failure Modes}\label{sec:decode:pitfalls}

\begin{pitfall}[KV Cache Memory Exhaustion \& Contiguous Fragmentation]
Pre-allocating contiguous KV buffers for $s_{\max}$ tokens per request wastes up to $60$--$80\%$ of GPU VRAM to internal and external fragmentation \bookcite{foundations_large_language_models}{Ch.~5, p.~237}. Under high concurrency, this triggers out-of-memory (OOM) evictions before compute is saturated. \emph{Remedy:} Use PagedAttention block tables, GQA/MQA, and chunked prefilling to cap per-iteration memory spikes.
\end{pitfall}

\begin{pitfall}[Logit Overflow, Temperature Collapse \& Top-$p$ Empty Support]
Dividing logits by tiny $\tau \to 0^+$ causes $\exp(u_j / \tau)$ to overflow FP16/BF16 ($> 65504$) into \texttt{NaN}. Conversely, applying top-$p$ without safeguarding the top-1 token can yield an empty set if numerical rounding drops $p_{(1)}$ below $p$, or if $-\infty$ masking is applied after softmax without renormalization \bookcite{build_large_language_model_from_scratch}{Sec.~5.3.3, pp.~394--395}. \emph{Remedy:} Route $\tau = 0$ directly to $\argmax$, apply log-sum-exp stabilization ($\vecb{u} - \max \vecb{u}$), and always retain at least $w_{(1)}$ in $\mathcal{V}_p$.
\end{pitfall}

\begin{pitfall}[Speculative Decoding Acceptance Collapse]
If the draft model $q$ diverges from $p$ (e.g., domain shift, mismatched tokenizer, or high temperature $\tau > 1$ amplifying tail differences), the expected acceptance rate $\alpha = \sum_v \min(p(v), q(v))$ drops sharply, causing wasted draft compute and slower-than-baseline latency. \emph{Remedy:} Dynamically adapt lookahead $\tau$ based on rolling acceptance rates or use self-speculative / Medusa-style heads aligned with $p$.
\end{pitfall}

\begin{pitfall}[Overthinking \& Verifier Misalignment in Test-Time Scaling]
In RLVR and PRM-guided search, unconstrained thinking budgets induce \emph{overthinking} (repetitive self-verification loops on trivial prompts that inflate latency and degrade accuracy) or reward hacking against imperfect PRMs \bookcite{reinforcement_learning_from_human_feedback}{Ch.~7, pp.~110--111}, \bookcite{foundations_large_language_models}{Ch.~5, p.~251}. \emph{Remedy:} Apply curriculum length penalties during RLVR training and calibrate PRM scores against ground-truth outcome verifiers.
\end{pitfall}

\end{document}
